\documentclass[12pt]{article}
\usepackage{amsmath,amssymb,mathtools,amsthm}
\usepackage{fullpage}
\usepackage[hidelinks]{hyperref}

\newcommand{\conv}{\operatorname{conv}}
\newcommand{\RR}{\mathcal R}
\newcommand{\BB}{\mathcal B}

\newtheorem{theorem}{Theorem}
\newtheorem{lemma}[theorem]{Lemma}
\newtheorem{proposition}[theorem]{Proposition}

% The supplied referee report records a warning-free, seven-page
% compilation of a copy with all three local corrections, with
% references and citations resolved.
%% FLAG: This session has no command-execution tool, so the requested
%% FLAG: pdflatex run could not be repeated here; the compilation
%% FLAG: results above are attributed to the referee, not to a new run.

\begin{document}

\begin{center}
\textbf{Crossing families via dual levels}
\end{center}

\noindent\textbf{Abstract.}
We prove that there is an absolute constant \(c>0\) such that every
set of \(n\geq2\) points in the plane with no three collinear contains
at least \(cn\) pairwise vertex-disjoint segments that cross pairwise
in their interiors. The proof combines an extraction theorem of
Pach, Rubin, and Tardos with a deletion argument in a dual arrangement
of lines. Controlling the deletion separately at each level avoids
a size-dependent loss.

\section{Introduction}\label{sec:introduction}

How large a family of pairwise crossing segments must a planar point
set contain? A point set is in \emph{general position} if no three of
its points are collinear. A \emph{crossing family} is a collection of
segments spanned by the point set whose endpoints are all distinct
and whose members cross pairwise in their relative interiors.
We prove a lower bound proportional to the number of points.

\begin{theorem}[Linear crossing families]\label{thm:main}
There is an absolute constant \(c>0\) such that every set of \(n\geq2\)
points in the plane in general position contains a crossing family
of size at least \(cn\).
\end{theorem}

The proof combines the near-avoidance extraction theorem of Pach,
Rubin, and Tardos \cite[Lemma~3.3]{PRT2021} with a deletion estimate
for two separated point sets. We first define the quantity that
controls the deletion.
% The external extraction result is Lemma 3.3 of the cited paper. ([arxiv.org](https://arxiv.org/pdf/1904.08845v2))

Two finite nonempty point sets \(A,B\) are \emph{separated} if their
convex hulls are disjoint. They can then be strictly separated by a
line. Now assume that \(A,B\) are separated, that \(A\cup B\) is
in general position, and that \(|A|=|B|=m\geq1\).
For distinct points \(p,q\), let \(\ell(p,q)\)
denote their entire supporting line. Define the \emph{avoidance defect}
by
\begin{align}
 t(A,B)={}&
 \bigl|\{\{a,a'\}\subseteq A:
     a\ne a',\ \ell(a,a')\cap\conv(B)\ne\varnothing\}\bigr|
 \notag\\
 &+\bigl|\{\{b,b'\}\subseteq B:
     b\ne b',\ \ell(b,b')\cap\conv(A)\ne\varnothing\}\bigr|.
 \label{eq:defect}
\end{align}
Each counted object is an unordered pair. General position ensures
that no point of the opposite set lies on the indicated line.
Consequently, the line meets the opposite convex hull precisely
when that set has points strictly on both sides of the line.

\begin{proposition}[Dual-level deletion bound]\label{prop:key}
Let \(m\geq1\) be an integer, and let \(A,B\) be separated
\(m\)-point sets in the plane whose union is in general position.
Put \(t=t(A,B)\). After deleting at most \(2t\) of the \(m^2\)
segments joining \(A\) to \(B\), the remaining segments can be
partitioned into \(2m-1\) possibly empty crossing families.
In particular, one such family has size at least
\begin{equation}
 \frac{m^2-2t}{2m-1}.
 \label{eq:keybound}
\end{equation}
If \(t\leq m^2/4\), there is a crossing family of size at least \(m/4\).
\end{proposition}

The proposition includes \(m=1\): the defect is then zero, and the
unique joining segment forms the required family.

The proof of Proposition~\ref{prop:key} groups joining segments by
the level of their corresponding dual intersections, where a level
records how many dual lines lie below an intersection. We mark
selected lines at intersections representing avoidance defects,
and delete joining segments whose dual intersections use marked
lines at the same level. The main difficulty is that a single
marked line may contain many such intersections.
Lemma~\ref{lem:repeat} charges repeated incidences to distinct
defects at that level, while Lemma~\ref{lem:clean} shows that the
retained intersections give a crossing family. Each defect
contributes to the deletion bound only at its own level, so summing
over levels introduces no additional factor depending on \(m\).

Section~\ref{sec:extraction} records the extraction theorem and its
specialization to complete geometric graphs.
Sections~\ref{sec:dual} and~\ref{sec:deletion} construct the dual
sweep and prove Proposition~\ref{prop:key}.
Section~\ref{sec:linear} chooses fixed parameters to deduce
Theorem~\ref{thm:main}.

\section{Near-avoidance extraction}\label{sec:extraction}

The following consequence of \cite[Lemma~3.3]{PRT2021} is the only
external mathematical ingredient. References to definitions and
lemmas in that paper use the numbering of arXiv:1904.08845v2.

\begin{theorem}[Near-avoidance extraction]\label{thm:extraction}
There is an absolute constant \(K\geq2\) such that, for every
\(0<\varepsilon<1\), every positive integer \(m\), and every
general-position point set \(P\) in the plane of cardinality
\[
 n\geq K\varepsilon^{-4}m,
\]
there are separated subsets \(A,B\subseteq P\) satisfying
\[
 |A|=|B|=m,
 \qquad
 t(A,B)\leq\varepsilon m^2.
\]
\end{theorem}
% The parameter dependence and quantifiers agree with Lemma 3.3
% in the cited arXiv version. ([arxiv.org](https://arxiv.org/pdf/1904.08845v2))

\begin{proof}[Derivation]
Write \(K_0\) for the absolute constant in
\cite[Lemma~3.3]{PRT2021}. That lemma applies to a geometric graph
with \(n\) vertices and at least \(\delta n^2\) edges when
\[
 n\geq \frac{K_0m}{\varepsilon^4\delta^5}.
\]
It provides separated \(m\)-point sets for which the total number
of incomparable pairs in the two avoidance orders is at most
\(\varepsilon m^2\). By \cite[Definition~3.1 and Lemma~2.2]{PRT2021},
these incomparable pairs are exactly the unordered pairs whose
supporting line meets the opposite convex hull. Their total is
therefore \(t(A,B)\).
% The identification of incomparability with the line--hull condition
% is stated in Lemma 2.2; Definition 3.1 sums the two counts. ([arxiv.org](https://arxiv.org/pdf/1904.08845v2))

For the complete geometric graph on \(n\geq2\) vertices,
\[
 \binom n2\geq \frac{n^2}{4}.
\]
Apply the lemma with \(\delta=1/4\), and take
\[
 K\geq\max\{2,4^5K_0\}.
\]
The hypothesis \(n\geq K\varepsilon^{-4}m\) then ensures both
\(n\geq2\) and the required size condition.
\end{proof}

\section{The dual arrangement}\label{sec:dual}

We now fix separated sets \(A,B\), each of size \(m\geq1\), whose
union is in general position. The aim is to represent each pair
counted by \(t(A,B)\) as a dual intersection at a single level.

Choose coordinates in which all points of \(A\) have negative first
coordinate and all points of \(B\) have positive first coordinate.
We may also assume that all first coordinates are distinct and that
distinct spanned lines have distinct slopes. To obtain these
additional conditions, make an arbitrarily small generic
perturbation. Because there are finitely many triples and every
triple orientation is nonzero, the perturbation can preserve all
orientation signs as well as strict separation. It therefore
preserves \(t(A,B)\) and the crossing relation between every pair
of spanned segments. Any crossing family obtained after perturbation
gives a crossing family with the same endpoint labels in the
original sets.

Dualize a point \(p=(p_x,p_y)\) to the line
\[
 p^*(x)=p_xx-p_y.
\]
The lines dual to \(A\) form a red family \(\RR\), all of negative
slope; those dual to \(B\) form a blue family \(\BB\), all of
positive slope. The distinct first coordinates give distinct dual
slopes, and general position excludes three concurrent dual lines.
Moreover, the abscissa of \(p^*\cap q^*\) is the slope of
\(\ell(p,q)\). Thus all dual intersection abscissas are distinct.

For \(k\in\{1,\ldots,2m-1\}\), a dual intersection has
\emph{level \(k\)} if exactly \(k-1\) dual lines lie strictly below
it. At an abscissa \(x\) containing no intersection, let \(L_k(x)\)
be the set of the \(k\) lowest dual lines, ordered by their values
at \(x\). As \(x\) increases through a level-\(k\) intersection,
the two incident lines exchange ranks \(k\) and \(k+1\): one
leaves \(L_k\), and the other enters. Intersections at other
levels leave \(L_k\) unchanged.

An intersection of a red and a blue line is \emph{bichromatic};
an intersection of two lines of the same color is
\emph{monochromatic}. Every red slope is smaller than every blue
slope. Hence, before a bichromatic intersection the red line lies
above the blue line, and afterward it lies below. At a
bichromatic level-\(k\) intersection, the red line therefore
enters \(L_k\) and the blue line leaves \(L_k\).

A red--red intersection is \emph{defective} if some blue line lies
strictly above it and some blue line lies strictly below it.
Define defective blue--blue intersections by interchanging the
colors. Let
\begin{align*}
 T_k^R&=\#\{\text{defective red--red intersections of level }k\},\\
 T_k^B&=\#\{\text{defective blue--blue intersections of level }k\},
\end{align*}
and put \(T_k=T_k^R+T_k^B\).

These intersections encode the avoidance defect exactly. For example,
if \((s,h)=a^*\cap a'^*\), then \(\ell(a,a')\) has equation
\(y=sx-h\). Blue dual lines lie both above and below \((s,h)\)
precisely when points of \(B\) lie on both sides of \(\ell(a,a')\).
The analogous statement holds for pairs in \(B\). Since every
dual intersection has a unique level,
\begin{equation}
 \sum_{k=1}^{2m-1}T_k=t(A,B).
 \label{eq:sumdefects}
\end{equation}

It remains to show that, at each level \(k\), deleting at most
\(2T_k\) bichromatic intersections leaves a crossing family in the
primal plane.

\section{Levelwise deletion}\label{sec:deletion}

Fix a level \(k\). All markings in this section refer only to this
level; marks do not carry from one level to another.

At each defective red--red level-\(k\) intersection, mark the red
line that \emph{leaves} \(L_k\). At each defective blue--blue
level-\(k\) intersection, mark the blue line that \emph{enters}
\(L_k\). Thus each defective intersection marks one line, though a
line may receive several marks. A bichromatic level-\(k\)
intersection is \emph{clean} if neither incident line is marked.

The first lemma explains this directional choice. A red line must
leave \(L_k\) between successive bichromatic entrances; a blue
line must enter between successive bichromatic exits. The
intermediate changes that account for these repetitions are
necessarily defective.

\begin{lemma}[Repeated incidences]\label{lem:repeat}
For a red line \(r\), let \(d_k(r)\) be the number of bichromatic
level-\(k\) intersections on \(r\), and let \(u_k(r)\) count the
defective red--red level-\(k\) intersections at which \(r\) leaves
\(L_k\). Then
\[
 d_k(r)\leq1+u_k(r).
\]
For a blue line \(b\), define \(d_k(b)\) analogously, and let
\(u_k(b)\) count the defective blue--blue level-\(k\) intersections
at which \(b\) enters \(L_k\). Then
\[
 d_k(b)\leq1+u_k(b).
\]
\end{lemma}

\begin{proof}
If \(d_k(r)\leq1\), the red inequality is immediate. Otherwise,
consider two consecutive bichromatic level-\(k\) intersections
on \(r\), with abscissas \(s<s'\) and blue partners \(b,b'\),
respectively. Immediately after \(s\), the line \(r\) belongs to
\(L_k\); immediately before \(s'\), it does not. It must therefore
leave \(L_k\) at an intermediate level-\(k\) intersection.
A red line can only enter, not leave, at a bichromatic
intersection, so this exit is red--red.

For every \(x\in(s,s')\), the slope inequalities give
\[
 b(x)>r(x)>b'(x).
\]
Indeed, \(b\) lies above \(r\) after their intersection at \(s\),
whereas \(b'\) lies below \(r\) before their intersection at \(s'\).
The red--red exit is therefore defective. The open intervals
between successive bichromatic intersections on \(r\) are
disjoint, so they require distinct defective exits. Thus
\(u_k(r)\geq d_k(r)-1\).

For a blue line \(b\), let two consecutive bichromatic
level-\(k\) intersections have abscissas \(s<s'\) and red
partners \(r,r'\). Immediately after \(s\), the line \(b\) is
outside \(L_k\); immediately before \(s'\), it is inside.
An intermediate entrance is therefore required. Since a blue
line leaves at a bichromatic intersection, this entrance is
blue--blue. Throughout \((s,s')\),
\[
 r(x)<b(x)<r'(x),
\]
so the entrance is defective. Again, disjoint intervals require
distinct entrances, proving the blue inequality.
\end{proof}

\begin{lemma}[Bounded deletion]\label{lem:deletion}
At most \(2T_k\) bichromatic level-\(k\) intersections are not clean.
\end{lemma}

\begin{proof}
Let \(S_R\) be the set of marked red lines. Each defective red--red
level-\(k\) intersection marks exactly its exiting line. Hence
\[
 |S_R|\leq T_k^R,
 \qquad
 \sum_{r\in S_R}u_k(r)=T_k^R.
\]
Lemma~\ref{lem:repeat} gives
\[
 \sum_{r\in S_R}d_k(r)
 \leq |S_R|+\sum_{r\in S_R}u_k(r)
 \leq2T_k^R.
\]
This bounds the number of bichromatic level-\(k\) intersections
incident to marked red lines.

For the set \(S_B\) of marked blue lines, the same count uses
entrances rather than exits:
\[
 |S_B|\leq T_k^B,
 \qquad
 \sum_{b\in S_B}u_k(b)=T_k^B,
 \qquad
 \sum_{b\in S_B}d_k(b)\leq2T_k^B.
\]
Every non-clean intersection is incident to a marked red line or
a marked blue line. Adding the two bounds, with possible double
counting, gives at most
\(2T_k^R+2T_k^B=2T_k\) non-clean intersections.
\end{proof}

\begin{lemma}[Clean intersections give a crossing family]
\label{lem:clean}
The primal segments corresponding to the clean bichromatic
intersections at level \(k\) have distinct endpoints and cross
pairwise in their interiors.
\end{lemma}

\begin{proof}
We first exclude shared endpoints. We then establish strict height
inequalities at each of two clean intersections; these inequalities
give the two side-of-line tests for a proper primal crossing.

An unmarked red line has \(u_k(r)=0\), so
Lemma~\ref{lem:repeat} shows that it is incident to at most one
bichromatic level-\(k\) intersection. The same conclusion holds
for an unmarked blue line. Thus clean intersections use distinct
dual lines, and the corresponding primal segments have distinct
endpoints. If there are fewer than two clean intersections, the
crossing assertion is vacuous.

Consider two clean intersections
\[
 q=(s,h)=r\cap b,
 \qquad
 q'=(s',h')=r'\cap b',
 \qquad s<s',
\]
where \(r,r'\) are red and \(b,b'\) are blue. We first prove
\begin{equation}
 r'(s)>h>b'(s).
 \label{eq:firstbracket}
\end{equation}
The four lines are distinct, and no three dual lines concur, so
neither comparison can be an equality.

Suppose that \(r'(s)<h\). Then \(r'\) belongs to \(L_k\) near
\(s\), but it is outside \(L_k\) immediately before its
bichromatic intersection at \(s'\). It must leave \(L_k\) at
an intermediate level-\(k\) intersection. This exit is red--red,
because a red line enters at a bichromatic intersection.

Throughout \((s,s')\), we have
\[
 r'(x)<b(x),
 \qquad
 r'(x)>b'(x).
\]
For the first inequality, use \(r'(s)<b(s)=h\) and the strictly
negative slope of \(r'-b\). For the second, use the fact that
\(r'\) lies above \(b'\) before their intersection at \(s'\).
The red--red exit is consequently defective and marks \(r'\),
contrary to the cleanliness of \(q'\). Hence \(r'(s)>h\).

Now suppose that \(b'(s)>h\). The line \(b'\) is outside \(L_k\)
near \(s\), but inside immediately before \(s'\). It must enter
\(L_k\) at an intermediate blue--blue level-\(k\) intersection.
Throughout \((s,s')\),
\[
 b'(x)>r(x),
 \qquad
 b'(x)<r'(x).
\]
The first inequality follows from \(b'(s)>r(s)\) and the strictly
positive slope of \(b'-r\). The second holds before the
intersection of \(b'\) and \(r'\). Thus this entrance is
defective and marks \(b'\), again contradicting the cleanliness
of \(q'\). This proves \eqref{eq:firstbracket}.

To obtain the inequalities at \(q'\), reverse the horizontal
coordinate and interchange the color names. Red lines again have
negative slopes and blue lines positive slopes. Vertical ranks
and defectiveness are unchanged. Under the reversed sweep, an old
red exit becomes a new blue entrance, and an old blue entrance
becomes a new red exit. The marking rule, and hence cleanliness,
is therefore preserved. Applying \eqref{eq:firstbracket} with the
two intersections in reversed order gives
\begin{equation}
 b(s')>h'>r(s').
 \label{eq:secondbracket}
\end{equation}

We now translate the two brackets back to the primal plane.
Let \(a\in A\) and \(c\in B\) correspond to \(r\) and \(b\).
Their supporting line has equation \(y=sx-h\), and for any
primal point \(p\),
\[
 \det(c-a,p-a)
 =-(c_x-a_x)\bigl(p^*(s)-h\bigr).
\]
Since \(c_x-a_x>0\), equation~\eqref{eq:firstbracket} places the
endpoints of the other primal segment strictly on opposite sides
of \(\ell(a,c)\). Equation~\eqref{eq:secondbracket} gives the
reciprocal statement. The supporting lines therefore meet in the
relative interior of both segments, as required.
\end{proof}

\begin{proof}[Proof of Proposition~\ref{prop:key}]
Each of the \(m^2\) joining segments corresponds to a bichromatic
intersection at one of the \(2m-1\) levels. At every level, discard
precisely the intersections that are not clean.
Lemma~\ref{lem:deletion} and \eqref{eq:sumdefects} bound the total
number discarded by
\[
 2\sum_{k=1}^{2m-1}T_k=2t.
\]
By Lemma~\ref{lem:clean}, the retained intersections at each level
give a crossing family. These families partition the retained
segments, since every dual intersection has a unique level.
The perturbation described in Section~\ref{sec:dual} preserves
both the defect and all crossing relations, so the same partition,
with the same endpoint labels, applies to the original sets.

Averaging the cardinalities of the \(2m-1\) families proves
\eqref{eq:keybound}. If \(t\leq m^2/4\), then
\[
 \frac{m^2-2t}{2m-1}
 \geq \frac{m^2}{2(2m-1)}
 \geq \frac m4. \qedhere
\]
\end{proof}

\section{Proof of the linear bound}\label{sec:linear}

We use the extraction theorem with the fixed error
\(\varepsilon=1/4\), for which Proposition~\ref{prop:key} retains
a crossing family of size at least one quarter of either extracted
set. It remains to choose the extraction scale and account for
small point sets.

\begin{proof}[Proof of Theorem~\ref{thm:main}]
Take \(K\) from Theorem~\ref{thm:extraction}, choose an integer
\[
 H\geq K4^4,
\]
and set \(c=1/(8H)>0\). These constants are absolute; in particular,
\(H\geq2\).

Let \(P\) be a general-position set of \(n\geq2\) points in the
plane. First suppose that \(n\geq2H\), and put
\(m=\lfloor n/H\rfloor\). Then
\[
 n\geq Hm\geq K4^4m,
 \qquad
 m\geq \frac{n}{2H}.
\]
Theorem~\ref{thm:extraction}, applied with \(\varepsilon=1/4\),
provides separated subsets \(A,B\subseteq P\) satisfying
\[
 |A|=|B|=m,
 \qquad
 t(A,B)\leq \frac{m^2}{4}.
\]
Proposition~\ref{prop:key} gives a crossing family of size at least
\[
 \frac m4\geq \frac{n}{8H}=cn.
\]
All its segments are spanned by the original point set \(P\).

If \(2\leq n<2H\), choose any two points of \(P\). Their segment
is a crossing family of size one, and
\[
 1\geq \frac{n}{8H}=cn.
\]
This covers all \(n\geq2\).
\end{proof}

\begin{thebibliography}{9}

\bibitem{PRT2021}
J.~Pach, N.~Rubin, and G.~Tardos.
\newblock Planar point sets determine many pairwise crossing segments.
\newblock \emph{Advances in Mathematics} \textbf{386} (2021), 107779.
\newblock doi: \nolinkurl{10.1016/j.aim.2021.107779}.
% The journal metadata and DOI agree with the publisher's record. ([sciencedirect.com](https://www.sciencedirect.com/science/article/pii/S0001870821002188?utm_source=openai))
\newblock Also arXiv:1904.08845v2, April 30, 2023.
% The arXiv submission history records this version and date. ([arxiv.org](https://arxiv.org/abs/1904.08845v2))

\end{thebibliography}

\end{document}
